\section{A Training-Free State-Conditioned Score}
\label{sec:analysis}

\subsection{Definition}

We use a state-conditioned closed form as a proposal signal, not as an unbiased estimator of the
true marginal gain.  Given the realised exposure state, let $g(e)=2e(1-e)$ and let
$N^+(w)$ be the out-neighbours of a candidate $w$, with edge weight $\omega_{wv}$.  Adding $w$
raises the exposure of $v$ by
\begin{equation}
  \Delta e_w(v) = \omega_{wv}\bigl(1-\delta(v)\bigr).
  \label{eq:delta-e}
\end{equation}
The one-hop proposal score is
\begin{equation}
  \mathrm{score}(w) = \sum_{v\in N^+(w)}
  \left[g\bigl(\delta(v)+\Delta e_w(v)\bigr)-g\bigl(\delta(v)\bigr)\right].
  \label{eq:score}
\end{equation}
We also use a two-hop version, denoted $\delta_2$, in which the exposure increment is propagated
once through the out-edges of each affected node, with the same $(1-\delta(\cdot))$ attenuation.

The score has two useful properties.  It is training-free and requires no candidate-level
Monte-Carlo trials after the exposure state is available.  It is also non-monotone: a candidate
that pushes an already exposed neighbour past its window can receive a negative score.  These
properties make it a natural proposal signal for the corrected regime sweep in
Section~\ref{sec:regime}.

\subsection{Scope of the score}

The closed form is not asserted to be an unbiased estimator, a certificate, or a replacement for
the true objective.  Its current evidence is limited to the MC=300 rank comparison in
Section~\ref{sec:regime}; all earlier calibration and policy-comparison tables were produced
before the state-machine correction or below the required Monte-Carlo budget and are omitted.
Whether $\delta_2$ can support a reproducible quality--cost advantage in sequential seed selection
must be established by a new post-fix experiment with Full-MC, degree, random-pruning, repeated
candidate pools, and explicit state-acquisition accounting.
